% ======================================================================
% DRAFT v1 — Stellar Mass Shapes Gravitational Bonding in Multi-Planet Systems
% Leslie Yarbrough
% April 2026
% ======================================================================

\documentclass[twocolumn]{aastex631}

\usepackage{graphicx}
\usepackage{amsmath}
\usepackage{amssymb}
\usepackage{hyperref}

\shorttitle{Stellar Mass Shapes Gravitational Bonding}
\shortauthors{Yarbrough}

\begin{document}

\title{Stellar Mass Shapes Gravitational Bonding in Multi-Planet Systems}

\author{Leslie Yarbrough}
\affiliation{Vela Research Collective}

\begin{abstract}
In multi-planet systems, adjacent planets often lock into simple-integer period
ratios---gravitational ``bonds.'' Analyzing 336 multi-planet systems (1182
planets) from the NASA Exoplanet Archive, we find that 56.6\% of adjacent planet
pairs are in exact resonance ($\pm 3\%$ of a small integer ratio). The fraction
of bonded pairs depends strongly on host star mass, and the relationship is
non-monotonic: M dwarfs ($M_\star < 0.6\,M_\odot$) show the highest bond
fraction (68\%), F dwarfs ($M_\star \geq 1.1\,M_\odot$) are close behind (61\%),
while G dwarfs ($0.9 \leq M_\star < 1.1\,M_\odot$) are lowest (50\%). A
stability-filtered null ensemble---randomized period architectures that satisfy
Hill stability---reveals that G-dwarf systems are indistinguishable from random
packing within stable bounds (observed 52.6\% vs.\ null 52.9\%). M and F dwarfs
show genuine excess bonding above the null ($+5.2\%$ and $+6.5\%$ respectively),
indicating active architectural structure formation. G-dwarf bonds are also the
loosest (largest deviation from exact resonance), contradicting the hypothesis
that only deeply captured resonances survive around Sun-like stars. We interpret
the ``U-shape'' as reflecting two distinct disk-driven mechanisms: M dwarfs
produce bonds through long-lived, slow migration (careful settling), while F
dwarfs achieve similar bond fractions through rapid, vigorous migration
(multiple capture opportunities). G dwarfs occupy an intermediate regime where
neither path is efficient, producing systems that are dynamically stable but
architecturally unstructured. The Solar System---orbiting a G dwarf at the
bottom of the bonding curve---is a noble-gas-like system: mostly non-bonded,
with one clean resonance (Neptune-Pluto 3:2) and two loose near-bonds. The
emerging picture suggests that gravitational ``chemistry'' is a specific
phenomenon of certain stellar environments, not a universal outcome of planet
formation.
\end{abstract}

\section{Introduction}
\label{sec:intro}

\epigraph{It was the best of stars, it was the worst of stars, it was the age
of bonding, it was the age of isolation\ldots}{with apologies to C. Dickens}

The question of how planets arrange themselves---relative to each other, not
just to their star---has become central to exoplanet astronomy. Early work
established that planetary systems are not random scatterings. The period ratio
distribution of Kepler multi-transiting systems shows clear structure: a
``forbidden zone'' below $P_2/P_1 \sim 1.2$, a pile-up near 1.5, and a deficit
at 2.0\citep{fabrycky2014,fang_margot_2012}. Systematic trends in planetary
size within systems---peas-in-a-pod---point to shared formation
conditions\citep{weiss2018,millholland2021}. The mutual Hill spacing
distribution clusters around $\Delta \sim 12$--$20$ for Kepler
systems\citep{puwu2015,fang_margot_2013}.

But one dimension remains largely unexplored: the host star itself. Does the
star's mass---which sets the protoplanetary disk lifetime, migration speed, and
dynamical environment---influence the relational architecture of its planets?

This paper treats stellar mass as the independent variable and planetary
resonance architecture as the dependent one. We adopt a chemical framing:
simple-integer period ratios (3:2, 2:1, 4:3, 5:3) correspond to stable resonant
configurations. We call these ``bonds'' deliberately---they are quantized,
stable states that adjacent planets settle into through dynamical evolution.
This is not a metaphor. It is a structural claim: gravitational interactions in
multi-planet systems produce discrete relational configurations, and the
conditions for forming those configurations depend on the host star in a
systematic, non-monotonic way.

\section{Data \& Methods}
\label{sec:methods}

\subsection{Dataset}

We use the NASA Exoplanet Archive \texttt{pscomppars} table, accessed 2026-04-27.
Our sample comprises 336 systems with $\geq 3$ confirmed planets, totaling 1182
planets. Discovery methods include transit (907 planets), radial velocity (248),
TTV (17), pulsar timing (3), eclipse timing (3), and direct imaging (4).

For each adjacent pair, sorted by orbital period, we compute the period ratio
$P_2/P_1$. Semimajor axes are derived from orbital periods via Kepler's third
law where not directly available.

\subsection{Resonance Classification}

We define an ``exact resonance'' (or ``bond'') as a period ratio within
$\pm 3\%$ of a small integer ratio. Classification uses the \textit{nearest}
resonance by fractional deviation:

\begin{equation}
  \delta = \min_{(p:q)}\left|\frac{P_2/P_1}{(p/q)} - 1\right|
\end{equation}

where $(p:q) \in \{3:2, 2:1, 4:3, 5:3, 5:4, 7:5, 8:5, 9:5, 7:3, 5:2, 8:3,
3:1, 7:2, 4:1\}$. If $\delta \leq 0.03$, the pair is classified as the
corresponding resonance type. ``Near resonance'' requires $0.03 < \delta \leq
0.10$.

A \textit{pure molecule} is a system where every adjacent pair is in exact
resonance. A \textit{resonance chain} is $\geq 3$ consecutive bonded pairs.

\subsection{Stellar Mass Bins}

Host stars are binned by mass:
\begin{itemize}
  \item M dwarf: $M_\star < 0.6\,M_\odot$
  \item K dwarf: $0.6 \leq M_\star < 0.9\,M_\odot$
  \item G dwarf: $0.9 \leq M_\star < 1.1\,M_\odot$
  \item F dwarf: $M_\star \geq 1.1\,M_\odot$
\end{itemize}

We exclude stars above $1.6\,M_\odot$ to avoid evolved and A-type hosts.

\subsection{Detection Bias Mitigation}

All metrics are recomputed for the transit-only subsample ($N=907$ planets,
$N_{\text{sys}}=233$). The transit-only results are used as the primary
comparison to control for RV selection effects. We also verify results with a
$\geq 4$ planet transit-only subsample.

\subsection{Bond Tightness Metric}

For bonded pairs, we measure the fractional deviation:

\begin{equation}
  \delta_{\text{bond}} = \frac{|P_2/P_1 - (p/q)|}{(p/q)}
\end{equation}

Lower $\delta_{\text{bond}}$ corresponds to tighter, deeper resonance locks.

\subsection{Stability-Filtered Null Ensemble}

To test whether observed bonding exceeds what would be expected from random
dynamical survival, we construct a null ensemble. For each real system, we:
(1) preserve host star mass, planet count, inner/outer orbital periods, and
planet masses; (2) randomize interior planet periods in log-period space;
(3) recompute semimajor axes from randomized periods; (4) reject systems with
adjacent Hill spacings below 8 mutual Hill radii (Gyr-scale stability
threshold); (5) measure bond fraction in the surviving randomized systems.
This produces, for each system, a distribution of bond fractions expected from
stable random packing. We then compare observed bond fractions to this null.

\section{Results}
\label{sec:results}

\subsection{Global Bond Statistics}

Of 846 adjacent planet pairs across all 336 systems, 479 (56.6\%) are in exact
resonance, 240 (28.4\%) are near resonance, and 127 (15.0\%) are non-bonded.
Ninety-five systems (28\%) are pure molecules; 42 systems (13\%) have resonance
chains of 3 or more.

The most common bonding types are: 3:2 (90 pairs, 18.8\%), 2:1 (60, 12.5\%),
5:3 (55, 11.5\%), 9:5 (47, 9.8\%), 7:3 (39, 8.1\%), 5:2 (35, 7.3\%), 8:5 (30,
6.3\%), 4:3 (29, 6.1\%), 7:5 (21, 4.4\%), 8:3 (19, 4.0\%), 3:1 (17, 3.5\%),
7:2 (13, 2.7\%), 4:1 (13, 2.7\%), 5:4 (11, 2.3\%).

\subsection{The U-Shape: Bond Fraction by Stellar Mass}

\begin{table}
\centering
\caption{Bond fraction and related metrics by stellar mass bin. ``Pure mol.''
shows the fraction of systems where every adjacent pair is in exact resonance.
``Chain'' is mean consecutive bonded pairs.}
\begin{tabular}{lcccc}
\hline
Star type & $N_{\text{sys}}$ & Bond frac. & Pure mol. & Mean chain \\
\hline
M dwarf    & 49  & 68\% & 37\% & 1.65 \\
K dwarf    & 122 & 57\% & 27\% & 1.39 \\
G dwarf    & 116 & 50\% & 23\% & 1.20 \\
F dwarf    & 49  & 61\% & 35\% & 1.41 \\
\hline
\end{tabular}
\label{tab:ushape}
\end{table}

Table~\ref{tab:ushape} shows the central result: bond fraction follows a
U-shaped dependence on stellar mass. M dwarfs have the highest bond fraction
(68\%), G dwarfs the lowest (50\%), and F dwarfs (61\%) are close to M dwarfs.
The same pattern holds for pure molecule fraction (M: 37\%, G: 23\%, F: 35\%)
and mean chain length (M: 1.65, G: 1.20, F: 1.41).

The U-shape survives every filter we apply. Under transit-only, $\geq 3$ planets:
M (71\%), K (60\%), G (53\%), F (65\%). Under transit-only, $\geq 4$ planets:
M (78\%), K (67\%), G (62\%), F (70\%). The pure molecule gap widens under
transit-only: M dwarfs 41\% vs.\ G dwarfs 25\%---a 16-point gap compared to 14
in the full sample. This rules out simple RV detection bias as the cause.

\subsection{Bond Quality: Tightness}

\begin{table}
\centering
\caption{Bond tightness: mean fractional deviation from exact resonance.
Lower values = tighter bonds = deeper resonance locks.}
\begin{tabular}{lcccc}
\hline
Bond type & M dwarf & K dwarf & G dwarf & F dwarf \\
\hline
3:2        & \textbf{1.07\%} & 1.37\% & 1.49\% & 1.46\% \\
2:1        & \textbf{1.76\%} & 1.37\% & 1.98\% & 2.09\% \\
4:3        & 1.16\% & \textbf{1.07\%} & 1.81\% & 0.64\% \\
5:3        & \textbf{1.31\%} & 1.54\% & 1.75\% & 1.31\% \\
\hline
\end{tabular}
\label{tab:tightness}
\end{table}

M dwarfs produce the tightest bonds (Table~\ref{tab:tightness}). G dwarfs
produce the loosest bonds in 5 of 7 major resonance types. This is the opposite
of the ``survivor bias'' prediction---if G dwarfs bond less, one might expect
the surviving bonds to be deeper---and supports a model where bond quality
reflects post-capture dynamical settling time, not selective survival.

\subsection{Stability-Filtered Null}

\begin{table}
\centering
\caption{Observed vs.\ stability-filtered null bond fractions. $N_{\text{stab}}$
is the number of systems with at least one stable randomized configuration.
$\Delta f$ is observed minus null.}
\begin{tabular}{lcccc}
\hline
Host bin & $N_{\text{sys}}$ & Observed & Null & $\Delta f$ \\
\hline
low-mass ($<$0.6)   & 47 & 68.6\% & 63.5\% & $+5.2\%$ \\
K-like (0.6--0.8)   & 62 & 58.2\% & 57.5\% & $+0.8\%$ \\
solar-mass (0.8--1.2) & 197 & 52.6\% & 52.9\% & $-0.3\%$ \\
high-mass (1.2--1.6)  & 30 & 60.5\% & 54.0\% & $+6.5\%$ \\
\hline
\end{tabular}
\label{tab:null}
\end{table}

The stability-filtered null (Table~\ref{tab:null}) reveals that the G-dwarf
bond fraction is indistinguishable from random packing within stability bounds
(52.6\% observed vs.\ 52.9\% null). M and high-mass dwarfs show genuine excess
bonding above the null. This means:

\begin{enumerate}
  \item G dwarfs are at the \textit{dynamical floor}: their resonance structure
        is what you would expect from random dynamical survival alone.
  \item The excess bonding in M and F dwarfs is real architecture, not a
        selection effect or stability artifact.
  \item Gravitational ``chemistry''---bonded structure beyond random
        survival---is a specific phenomenon of certain stellar environments.
\end{enumerate}

\subsection{Notable Systems}

\textbf{TRAPPIST-1} (M dwarf, 0.09\,$M_\odot$): 7 planets, chain of 6
consecutive resonances. A heteropolymer:
8:5$-$5:3$-$3:2$-$3:2$-$4:3$-$3:2.

\textbf{HD~110067} (K dwarf, 0.80\,$M_\odot$): 6 planets, chain of 5.
A near-uniform crystal:
3:2$-$3:2$-$3:2$-$4:3$-$4:3.

\textbf{TOI-1136} (G dwarf): 6 planets, mixed chain with 2:1, 7:5, and 3:2s.

\textbf{Kepler-223} (G dwarf): 4 planets, textbook 4:3$-$3:2$-$4:3 chain---a
rare G-dwarf molecule that supports the rule.

\section{Two Mechanisms, One U-Shape}
\label{sec:mechanisms}

The U-shape---high bonding at both ends of the mass spectrum, low in the
middle---demands an explanation. The data rule out simple detection bias.
Disk evolution models offer a coherent framework.

\subsection{M Dwarf Mechanism: Slow Cooking}

M dwarfs have the longest-lived protoplanetary disks (5--10~Myr, compared to
2--5~Myr for G dwarfs)\citep{yasui2010,haisch2001}. Combined with low disk
masses, this produces slow planetary migration. Planets approach resonances
gradually, giving high capture probability. The extended disk lifetime also
provides post-capture dynamical settling: millions of years for orbits to settle
deeper into resonance.

This predicts: many bonds, tight bonds, diverse bond types. The data confirm
all three. M dwarfs have the highest bond fraction, the tightest bonds, and the
widest bond palette (2:1, 5:3, 7:3, and 8:3 all prominent). The chemistry
analog is low-temperature, long-duration reaction chemistry: thermodynamic
products, many stable outcomes.

\subsection{F Dwarf Mechanism: Fast Chemistry}

F dwarfs have short-lived (1--3~Myr) but massive disks\citep{sicilia_aguilar}.
Migration is fast and vigorous. Planets sweep through multiple resonances
rapidly, giving capture opportunities through a different channel: kinetic
capture rather than slow settling. The short disk lifetime leaves less time for
post-capture settling, so bonds are looser on average.

This predicts: comparable bond fraction to M dwarfs, but wider deviations. The
data confirm this: F dwarfs have 61\% bond fraction (close to M dwarfs' 68\%),
but their bonds are consistently looser (2.09\% deviation for 2:1 vs.\ M dwarfs'
1.76\%; 1.46\% vs.\ 1.07\% for 3:2). The chemistry analog is
high-temperature, short-duration reaction chemistry: kinetic products, many but
loose.

\subsection{G Dwarf Mechanism: The Dynamic Floor}

G dwarfs fall in the intermediate regime: disk lifetimes of 2--5~Myr, moderate
disk masses, moderate migration speeds. Neither slow enough for careful M-dwarf
locking nor fast enough for F-dwarf kinetic capture. Resonant encounters are
frequent but rarely successful: most are ``near misses'' where planets approach
a resonance but cross through it before capture can complete.

The data support this picture. G dwarfs show the lowest bond fraction, the
loosest bonds, and the narrowest bond palette. And crucially, the
stability-filtered null shows that G-dwarf systems are indistinguishable from
random packing within stability bounds. They are not building structure; they
are merely surviving.

The chemistry analog is the ``dead zone'' in reaction yield vs.\ temperature.
Neither too cold nor too hot: too moderate to be productive.

\subsection{The Deeper Pattern}

The U-shape may be a general feature of systems where bond formation requires
specific environmental conditions. Real chemistry shows a U-shaped reaction
yield with temperature---too cold and reactions don't proceed; too hot and
bonds break as fast as they form. The sweet spot is intermediate.

Planetary systems show the same pattern. M and F dwarfs both achieve bonding,
but via different mechanisms (slow time vs.\ high energy density). G dwarfs are
the lukewarm middle---neither cold-and-patient nor hot-and-aggressive. The
Solar System sits at this point: bonded in places (Neptune-Pluto 3:2) but
mostly unstructured.

This may be the deepest claim from this data: \textit{gravitational bonding
chemistry has its own ``reaction temperature,'' set by the host star's disk
properties.}

\section{The Solar System in Context}

Our Solar System orbits a G dwarf ($M_\odot = 1.0$). Table~\ref{tab:solarsys}
summarizes its planetary bonds.

\begin{table}
\centering
\caption{Adjacent planet pairs in the Solar System and their bond status.
``Tightness'' is the fractional deviation from the nearest integer ratio.}
\begin{tabular}{lccc}
\hline
Pair & Period ratio & Nearest resonance & Tightness \\
\hline
Mercury--Venus & 1.59 & 8:5 & 0.4\% (near) \\
Venus--Earth   & 1.63 & 5:3 & 2.5\% (near) \\
Earth--Mars    & 1.88 & 15:8 & --- \\
Mars--Jupiter  & 3.41 & 7:2 & 2.5\% (near)* \\
Jupiter--Saturn & 2.48 & 5:2 & 0.8\% (loose) \\
Saturn--Uranus & 2.85 & 20:7 & --- \\
Uranus--Neptune & 1.88 & 15:8 & --- \\
Neptune--Pluto & 1.50 & \textbf{3:2} & 0.01\% (clean) \\
\hline
\end{tabular}
\label{tab:solarsys}
\end{table}

*Mars--Jupiter: period ratio computed as Jupiter's period divided by Mars', but
the asteroid belt occupies the space between them, raising questions about
``adjacency'' in our framework.

The Solar System has one clean bond (Neptune--Pluto 3:2, the tightest
resonance in the entire dataset) and a few loose near-bonds. This is a typical
G-dwarf system: mostly non-bonded, occasional resonances, nothing tighter than
$\sim 1\%$ deviation.

The implication is worth stating clearly: \textbf{the Solar System is not a
typical outcome of planet formation. It is a specific outcome of forming around
a star with poor bonding conditions.} Surveys that extrapolate from ``the
Solar System is normal'' may systematically underestimate how structured
planetary systems can be.

\section{Conclusions}
\label{sec:conclusions}

\begin{enumerate}
  \item Of 846 adjacent planet pairs in 336 multi-planet systems, 56.6\% are in
        exact resonance ($\pm 3\%$ of a simple integer ratio). Gravitational
        bonding is the norm, not the exception.
  \item Bond fraction follows a U-shaped dependence on stellar mass. M dwarfs
        (68\%) and F dwarfs (61\%) bond strongly; G dwarfs (50\%) bond poorly.
  \item G dwarfs are at the dynamical floor: their bond fraction is
        indistinguishable from random packing within stable bounds. M and F
        dwarfs show genuine excess bonding.
  \item M dwarfs produce the tightest bonds (closest to integer ratios),
        consistent with long post-capture settling times. G dwarfs produce the
        loosest bonds.
  \item The U-shape is consistent with two distinct mechanisms: M-dwarf
        slow-capture (long disk lifetime, careful settling) and F-dwarf
        fast-capture (short disk lifetime, kinetic capture).
  \item The Solar System, orbiting a G dwarf at the bottom of the bonding
        curve, is a typical G-dwarf system: mostly unstructured, with one
        clean resonance and a few loose near-bonds.
\end{enumerate}

\subsection{Predictions}
\begin{itemize}
  \item Bond fraction should correlate with stellar age within each mass bin
        (older M dwarfs $\rightarrow$ more settling $\rightarrow$ tighter bonds).
  \item High-resolution ALMA observations of M-dwarf disks should confirm
        longer dissipation timescales.
  \item Bond fraction around F dwarfs should anticorrelate with disk accretion
        rate (faster accretion $\rightarrow$ less settling).
  \item Direct imaging surveys of M-dwarf systems should preferentially find
        resonance chain architectures.
\end{itemize}

\subsection{Future Work}
\begin{itemize}
  \item Population synthesis with disk evolution models to test the
        two-mechanism hypothesis quantitatively.
  \item N-body integrations comparing stability lifetimes of M-dwarf vs.
        G-dwarf resonance chains.
  \item Extension to non-adjacent pairs (1,3-interactions)---the gravitational
        analog of ring strain in molecular chemistry.
  \item Application of bond classification to the full Kepler DR25 catalog.
\end{itemize}

\bibliography{references}
\bibliographystyle{aasjournal}

\end{document}
